\subsection{ERET Routing and Atomicity}
ERET classifies its source before any stack access. (:ref:base.registers.STATE.STATUS.PM:)\texttt{=1}, (:ref:base.registers.STATE.STATUS.EA:)\texttt{=0}, (:ref:base.registers.STATE.STATUS.EDEPTH:)\texttt{=0}, (:ref:base.registers.STATE.STATUS.UO:)\texttt{=0}, and (:ref:base.control_registers.UCTL.V:)\texttt{=1} performs direct initial user entry from the bank. (:ref:base.registers.STATE.STATUS.PM:)\texttt{=1}, (:ref:base.registers.STATE.STATUS.EA:)\texttt{=1}, (:ref:base.registers.STATE.STATUS.EDEPTH:)\texttt{=1}, (:ref:base.registers.STATE.STATUS.UO:)\texttt{=1}, and (:ref:base.control_registers.UCTL.V:)\texttt{=1} returns the outer user event from the bank. Other supervisor states with (:ref:base.registers.STATE.STATUS.EA:)\texttt{=1} and nonzero (:ref:base.registers.STATE.STATUS.EDEPTH:) restore a full frame from the current SS:SP. Every other combination raises (:ref:base.events.EXCEPTION.INVALID_CONTROL_TRANSITION:).

A U-bank return validates UCTL, UINFO, the user segment images, USP, and the executable UPC before committing. It restores user execution, clears (:ref:base.control_registers.UCTL.V:), (:ref:base.registers.STATE.STATUS.EDEPTH:), (:ref:base.registers.STATE.STATUS.UO:), (:ref:base.registers.STATE.STATUS.EA:), and current DFA, and never reads the stack. A stacked return validates the entire frame, requires saved (:ref:base.registers.STATE.STATUS.EDEPTH:) plus one to equal the current depth and saved (:ref:base.registers.STATE.STATUS.UO:) to match the current chain, then restores (:ref:base.events.frame.EVENT_FRAME.FRAME_CONTROL.STATUS:) and (:ref:base.events.frame.EVENT_FRAME.FRAME_CONTROL.SAVED_DFA:) atomically. A failed validation or target probe leaves the current context and return source unchanged.

Normative instruction-local behavior is defined by (:ref:base.instructions.SYSCALL:) and (:ref:base.instructions.ERET:).\par
